\chapter*{Preface}
\addcontentsline{toc}{chapter}{Preface}
\markboth{Preface}{Preface}

Every trial is an exercise in reasoning under uncertainty. A jury must decide whether a man it has never met killed a woman it never knew, on the strength of a few fibres, a smear of DNA, a confused witness and a story told by counsel. A judge in the Commercial Court must decide whether an insurer would have written a policy if it had been told what the broker knew. A district judge in Uttar Pradesh, in a state with millions of pending matters, must decide which ones to hear this week. None of these decisions is certain, and all of them have consequences. Probability is the language for the first fact and utility is the language for the second. Bayes' rule and the principle of maximising expected utility are the grammar that joins them.

This book argues that law is already Bayesian in its bones, even where it refuses to say so, and that saying so out loud makes legal reasoning clearer, fairer and cheaper. The argument is not new. Laplace and Condorcet wrote about juries, Poisson built a probability theory of judgments, and in the twentieth century Dennis Lindley, Stephen Fienberg, Jay Kadane, I. J. Good, Philip Dawid, David Schum, Colin Aitken and Franco Taroni showed how far the framework can be pushed in real cases. What is new is the setting. Machines now read pleadings, draft contracts and predict the outcomes of appeals. Formal proof assistants such as Lean can check mathematical arguments line by line, and legal scholars have begun to encode statutes in languages a computer can verify. Lindley argued that a court should decide by maximising expected utility, and whether a machine might one day combine probabilities and utilities coherently on its behalf is now a practical question. We think the answer depends on whether the machine is Bayesian.

The book grew out of a series of papers. Our paper on Bayesian interpretations of English tort law treats the doctrines of negligence, causation and misrepresentation as constraints on belief and on action, and uses the collapse of Greensill Capital as a worked example. The paper on Bayesian stories, written with Shiva Maharaj, gives a weight of evidence theory of narratives, using the murder of Meredith Kercher in Perugia as its case. The paper on the Pleadings Game builds deep generative Bayesian computation on Thomas Gordon's formal model of civil pleading. Our work on deep learning and large language models in law asks what the new architectures can and cannot automate. Those papers are the spine of Parts II and III, and the work on machines informs the closing chapters of Part VII.

Around that spine we have placed the work of others whom we admire. Stephen Fienberg's panel for the National Research Council, reported in \emph{The Evolving Role of Statistical Assessments as Evidence in the Courts}, remains one of the best accounts of how statistics actually enters litigation. Jay Kadane's work on whether electronic draw poker is a game of skill, on jury selection, and with David Schum on the Sacco and Vanzetti evidence shows a Bayesian statistician doing the hard work of an expert witness. Part IV, Statisticians in the Courtroom, opens with them and then follows statistical evidence through five applied chapters on discrimination, securities fraud, medical murder cases, intellectual property and digital evidence.

Part V turns to comparative law. The common law of England and the civil law of France reason about evidence in different ways, and the legal origins literature of La Porta, Lopez de Silanes, Shleifer and Vishny famously argued that legal origin shapes investor protection and financial development. Africa, where the colonial powers left both systems side by side, is a natural laboratory, and Cameroon, which contains both within one state, is the sharpest experiment of all. India, with more than five crore pending cases, shows what happens when a common law system is overwhelmed by volume, and why Bayesian triage is a humane response.

Part VI follows Roger Gould into nineteenth century Corsica, where vengeance was a calibrated sanction rather than a blind passion, and Diego Gambetta into Sicily, where the mafia sold protection in a market that the state could not supply. These are societies that built order without courts, and they show why reputations and signals are, at bottom, posterior probabilities.

Part VII brings the economics of Gary Becker and Kevin Murphy to bear on costs, deterrence and illegal markets, explains how Lean style verification could make parts of legal decision making checkable, returns to Lindley's challenge about machines and the explainability of the algorithms that would answer it, and closes with a program and an outlook.

Our first goal is pedagogical, as stated at the front of the book. The book is organised as a casebook to that end, with leading cases in most chapters and Notes and Questions for class.

We have also tried to write for three audiences at once: lawyers who want to understand what a likelihood ratio is and why it matters, statisticians who want to see how their tools fit the institutions of the law, and policy makers who must decide how far to let algorithms into courtrooms. The mathematics is kept to what is needed, and each substantive chapter ends with Notes and Questions and a section of further reading. Where we describe live litigation we have relied on public judgments and official statements, and we have tried to be careful about what is established and what is alleged.

We thank Shiva Maharaj, our coauthor on the stories paper, and the many students at Chicago Booth and George Mason who have argued with us about juries, insurers and machines. Any errors are ours.

\medskip
\noindent Nick Polson, Chicago\\
Vadim Sokolov, Fairfax\\
November 2026
